\chapter*{Conferencia IX}
\addcontentsline{toc}{chapter}{LECTURA I - El timbre y el carácter del violín}
SEÑORES: En este momento se retoma la presentación de los modificadores de tono del violín, que actúan para disminuir el volumen y la intensidad del sonido. Al finalizar la sesión anterior, aún estábamos considerando la gran influencia que tiene la curvatura de las tapas sobre el tono. A continuación, presento este violín como otro ejemplo de curvatura, que invariablemente reduce la potencia sonora. Al observar este violín de perfil, se aprecia que toda la elevación de la curvatura se concentra en los primeros cinco centímetros de cada extremo de la tapa; que entre estas elevaciones abruptas, las tapas forman una línea recta; y que las demás dimensiones permanecen como de costumbre. La aplicación del arco demuestra que la potencia sonora de este violín, si bien es perceptiblemente mayor que la de un violín con una curvatura más pronunciada, aún carece de gran parte de su máximo potencial. En mi opinión, este modelo de curvatura para las tapas del violín ofrece una potencia sonora apenas ligeramente superior a la de un violín sin curvatura. De no ser por la ligera concentración de ondas sonoras en las salidas debido al arco lateral o transversal, no habría diferencia alguna en la potencia del sonido. Por lo tanto, la mayor disminución en el volumen e intensidad del sonido debido al arco, manteniendo las demás dimensiones iguales, se encuentra en los extremos; es decir, en la placa perfectamente plana y en el enorme arco mencionado anteriormente. Al exponer mi opinión sobre la causa de dicha disminución de la potencia del sonido, reconozco mi incapacidad para aportar pruebas. Si existiera un solo arco en las placas del violín, la prueba sería evidente.
150

en tan solo unos instantes. Es la presencia del arco lateral lo que añade complicación a este problema. Con solo el arco mayor (de extremo a extremo de las placas) a la vista, la línea de viaje seguida por las ondas sonoras, originadas en la caja de resonancia, podría trazarse con precisión en papel. Las leyes físicas que explican los movimientos de las ondas sonoras son fácilmente
comprendido. > ;
Las dos leyes relativas a este problema son: (a) Las ondas sonoras, en el momento de su origen, viajan en ángulo recto desde el agente que las produce; (b) Las ondas sonoras, tras incidir sobre una superficie reflectante, viajan desde ella en un ángulo igual al ángulo de incidencia. Es evidente que resulta sencillo trazar en papel una línea perpendicular a cualquier punto de la superficie interior de cualquier arco. También es evidente que una línea de movimiento de ondas sonoras, originada en cualquier punto de la caja de resonancia arqueada del violín, no incidirá en la parte posterior en un punto perpendicular a su origen, sino en puntos más cercanos a las salidas. Hasta aquí todo bien; pero el siguiente movimiento de la onda sonora es donde reside la dificultad. Es ahora cuando el arco lateral se hace presente. Es evidente que, al incidir sobre la parte posterior, la onda se reflejará; pero ¿en qué dirección viajará la reflexión? La función de los arcos es dirigir y concentrar el movimiento de las ondas en las salidas. También es la intención colocar las salidas en la línea de movimiento de las olas. Es evidente que tales intenciones se ven frustradas por enormes grados de arqueamiento lateral,
y también sin arquearse en absoluto. 151	Es evidente que

Cuando las placas carecen de curvatura, no puede haber una progresión directa de la onda sonora hacia las salidas. Por lo tanto, el sonido de dichas placas se ve disminuido en potencia. Es evidente que las líneas de propagación de las ondas sonoras dentro del violín de arco alto se desvían considerablemente de las salidas. En este caso, la pregunta es: "¿Dónde deberían ubicarse las salidas?".
[En la medida en que el violín no es más que un producto de la experiencia, por lo tanto podemos seguir esperando que la mejora del tono del violín solo siga a la experiencia. Con esta idea en mente, recuerdo una experiencia relacionada con un barril de sidra nueva que, sin querer, se dejó al sol abrasador durante todo el día. Siendo empleado como inspector de este barril en particular, y estando el barril bien cerrado, dirigí un golpe moderado con un martillo hacia las proximidades del tapón; entonces, y sin previo aviso, ese tapón salió volando hacia mi cara.
ftow (con ánimo vengativo) Sugiero hacer un agujero en la tapa armónica entre los pies del puente como lo apropiado para todos los violines que tienen una cintura estrecha.]
(5) El control de la acción de la caja de resonancia mediante la barra es fácil de lograr, muy fácil; me dan ganas de decir: "Demasiado fácil". La experiencia al lidiar con el problema de la barra del violín puede causar la misma decepción que la experiencia con cualquier otra "barra". En cualquier caso, "cargar" demasiado seguramente resultará decepcionante. La barra del violín es un factor poderoso en
152

Modificación del tono de las cuerdas G y D. La barra puede ser demasiado larga, gruesa o profunda, con una densidad de fibra excesiva y una masa demasiado ligera. El tono de las cuerdas G y D, incluso en la mejor caja de resonancia, puede verse fácilmente afectado por cualquiera de estos defectos, y arruinado por completo por su combinación. La posición de la barra también puede perjudicar seriamente el tono de las cuerdas G y D.
[Con el fin de que mis afirmaciones estén respaldadas por pruebas y de que puedan tener algún efecto, presento el caso de un violín que sufre una pérdida de timbre debido a un molesto "lobo" en el sol al aire, causado por la mala posición de la barra. Este violín está hecho de madera preciosa y, en todos los demás aspectos, su timbre es excelente. Su constructor es A.F. Anderson. En mi opinión, el trabajo mecánico de este violín es insuperable. El barniz es de la variedad más resistente y elástica. En mi opinión, a juzgar por la mala posición de la barra y la graduación de las tapas, el Sr. Anderson construye violines siguiendo reglas estrictas. Dado que el Sr. Anderson es un artesano de primera categoría, se espera que se complazca al señalar un error en su trabajo, por lo demás impecable.]
Los detalles se basan en la afirmación de que me opuse a abrir este violín. Mi reticencia se debía a la incertidumbre sobre la causa de este tono "lobo". De ninguna manera podía asumir una solución en este caso. Nunca antes me había encontrado con un tono lobo en la nota fundamental de la cuerda de sol como una falla de tono aislada. Muchas veces había logrado eliminar una "manada de lobos" del violín.
153

tono, pero "lobo" en un tono aislado es una propuesta diferente para la solución. Una vez me pidieron que eliminara "lobo" de la cuerda E de un violín viejo. Era un caso de "lobo" en sol y sol sostenido en alt.
Este viejo violín mostraba amplia evidencia de que generaciones habían dedicado su tiempo a mover el alma. En este trabajo, el borde derecho de la salida derecha se había desgastado una pulgada; y ambas placas estaban magulladas en un círculo de una pulgada de diámetro por los extremos móviles del alma. Al no poder determinar la causa del "tono lobo" con el arco, retiré la caja de resonancia. Aparentemente, la mano de obra dentro de este viejo violín era impecable. No pude ver ninguna causa del "tono lobo"; sin embargo, sabía con certeza que existía, y, como se indicó anteriormente. Con curiosidad por conocer la graduación de esta caja de resonancia, utilicé el calibrador y así encontré razones para "adivinar" la causa del tono "lobo" en este caso. Esto es lo que encontré con el calibrador: el constructor, ejerciendo una precaución extrema al reducir el grosor de la madera debajo de la cuerda Mi, había sobrepasado el límite de seguridad en tal trabajo; pero, solo un poco más allá. Bajo condiciones tan meteóricas como las que hasta ahora se han descrito cuidadosa y repetidamente, la intensidad del tono perteneciente a la cuerda Mi de este violín supera a todas las demás cuerdas Mi que he sometido a la prueba al aire libre. El tono de esta cuerda Mi se propagó a gran distancia de
1480 pies lineales.
[Tenga en cuenta esas condiciones meteóricas, porque un cambio en el único elemento de mayor cantidad de
154

El vapor de agua presente en el momento de realizar la prueba de larga distancia al aire libre podría duplicar, e incluso cuadruplicar, esta distancia.
La mayor potencia tonal que posee esta cuerda Mi concentra el interés en dos factores.
(a) Gradación de la caja de resonancia debajo de la cuerda Mi.
(b) Ampliación del área de la salida derecha. Brevemente, describiré la graduación debajo de esta cuerda E: En la posición del puente, el grosor es igual a 7,64 pulgadas; luego disminuye gradualmente hasta 4,64 en un punto a mitad de camino desde el puente hasta el extremo superior de la placa; desde este punto hasta el final de la placa, el grosor aumenta gradualmente hasta 9,64. Evidentemente, esta forma de caja de resonancia da como resultado la producción de dos resortes cónicos entre el puente y el extremo superior de la placa. Prácticamente, estos resortes cónicos están unidos a
el uno al otro en sus extremidades más delgadas.
Es evidente que el resorte más alejado del puente es el más fuerte. Es bien sabido que el resorte más fuerte actúa con mayor rapidez que el más débil. Es evidente que la acción de cualquiera de estos dos resortes debe excitar la acción del otro. Debido a los diferentes grados de grosor, es evidente que dicha acción debe variar en número por segundo; por lo tanto, estos dos resortes, al golpear el aire contenido con diferente número de golpes por segundo, producen dos tonos simultáneos; y, estos dos tonos, al estar afinados en tonalidades disonantes, operan para producir "lobo". Hay dos métodos
disponible para remedio:	(a). 155	Igualación	primavera

acción reduciendo el espesor de la más fuerte
resorte, (b). Igualar la acción del resorte acortando el resorte más débil. Debido al peligro de debilitar la potencia del tono con el primer método, por lo tanto elegí el segundo. Si el segundo
Si el primer método hubiera resultado infructuoso, habría comprendido la necesidad de emplearlo. Sin embargo, el segundo método resultó exitoso. Su aplicación práctica consistió en mover gradualmente tanto el puente como el alma hacia arriba o hacia adelante. Un fenómeno curioso apareció en el tono de la cuerda Mi a medida que el puente y el alma se acercaban a su posición final. Este fenómeno se manifestó mediante el desplazamiento del "lobo" de sol y sol sostenido a sol sostenido y la, luego a la y si bemol, y, en el siguiente movimiento, el "lobo" desapareció.
[La ampliación de la salida derecha en este violín, al no ser relevante, no requiere mayor atención por el momento.]
Vuelvo al violín de Anderson.
Al abrir este violín, la causa del "lobo" no era evidente a primera vista. Sin utilizar el calibrador, pude determinar que la graduación de la caja de resonancia no era la causa; sin embargo, en la caja de resonancia debe encontrarse la causa. Mientras observaba la superficie interior, y sintiéndome bastante inseguro de tener éxito en este caso, noté una oblicuidad inusual en la posición de la
bar.
[A veces el diagnóstico del médico de una enfermedad oscura debe depender de la negación; no es esto,
ni eso, ni el otro, y así sucesivamente por exclusión.
156

hasta que queden pocas, o una sola causa probable. Pero para el médico, hay una diferencia entre el paciente artificial y el paciente humano. Ante el primero, el médico no necesita molestarse en "parecer sabio".
Mediante mediciones, descubrí que el extremo superior de esta barra cruzaba la línea de la cuerda G, cruzaba la línea de la cuerda D y llegaba a un punto entre las cuerdas D y A. Dado que no encontré ningún otro detalle que pudiera explicar el efecto de "lobo" en la cuerda G al aire, retiré esta barra y coloqué otra con el centro de su grosor directamente debajo del centro del vástago izquierdo del puente, y con una oblicuidad igual a la de la cuerda G, ni mayor ni menor.
Cuando se volvió a reproducir el orden, no había ningún "lobo" en ninguna cuerda.
Al intentar explicar este caso de "lobo", no puedo presentar más que una opinión. Por lo tanto: creyendo que ciertas fibras de la caja de resonancia, debajo de cada cuerda, actúan para producir todos los tonos posibles en cada cuerda, por lo tanto, en este caso, creo que la gran oblicuidad de la barra causó
"lobo" sobre el tono G abierto, y mi razonamiento
es como sigue: La graduación de esta caja de resonancia, siendo una modificación de Stainer, por lo tanto, el mayor grosor estaba en la posición del puente, y desde allí el grosor disminuía hasta 4-64 en todos los puntos cercanos a los bordes y en los extremos de la placa. El grosor en el puente era igual a 10-64. Por lo tanto, la parte más delgada de esta caja de resonancia, debajo de la cuerda G, estaba encima de la barra; no encima
157

El extremo superior de la barra, el extremo superior de la barra se encontraba en un punto entre las cuerdas D y A. Teniendo en cuenta la gran oblicuidad de esta barra, es evidente que las fibras de la caja de resonancia debajo de las cuerdas G y D se acortan de izquierda a derecha. Así, las fibras más largas están debajo de la G, las fibras más cortas están debajo de un punto entre D y A. Es evidente que cuando estas longitudes variables de fibras de la caja de resonancia se activan produciendo sonido, debe haber más de un sonido, o un tono. No sé cómo determinar el número exacto de estos sonidos. Supongo que algunos de ellos, al estar en un tono inarmónico con el G abierto, operaban para producir "ruido" o "silbido de lobo".
tono a partir de ahí.
La curación fue completa.
En mi experiencia, aumentar la longitud de la barra más allá de 10 pulgadas reduce la potencia tonal y también la duración del sonido. Esta afirmación se hace entendiendo que la cantidad de madera tanto en la barra como en la caja de resonancia se ha ajustado con precisión para producir la máxima potencia tonal. Teniendo esto en cuenta, afirmo que acortar la barra, manteniendo las demás dimensiones iguales, reduce el tono, aumenta el volumen y disminuye la intensidad del sonido; asimismo, en las condiciones anteriores, aumentar el grosor o la profundidad de la barra eleva el tono y disminuye la potencia tonal. La posición de la barra puede aumentar la potencia tonal de la cuerda G y, al mismo tiempo, disminuir la potencia tonal de la cuerda D. Por lo tanto, cuando
158

La distancia entre el extremo superior de las salidas es igual a 1 y 5-8 o 1 y 11-16, entonces, colocar el lado izquierdo de la barra al ras con el extremo superior de la salida funciona para aumentar la potencia de la cuerda G mientras disminuye la potencia de la cuerda D.
